% Page style shared by single pages (page-master.tex) and the one-run book
% (build-book.mjs -> letteris-book.tex): everything between \documentclass and
% \begin{document}. The page's own settings (\pagecfg, \pageparity, ... from
% pageN.tex) are read when each page is shipped out, so one copy serves every
% page of the book. Per-page state is (re)set in page-content.tex.

\usepackage{ifthen}
% Common preamble for all retypeset pages - geometry, fonts, base macros.
% Each page includes this, then calls \setuppage{} with its chapter/catchword data.

\usepackage[paperwidth=7.956in,paperheight=14.322in,
            textwidth=6.655in,marginparwidth=0.36in,marginparsep=0.1in,
            top=0.96in,bottom=1.56in,
            headheight=15pt,headsep=6mm,footskip=6mm]{geometry}
\usepackage{fontspec}
\usepackage{polyglossia}
\usepackage{subscript}
\usepackage{fancyhdr}
\usepackage{tikz} % \cutletter (clipping)
\setmainlanguage{hebrew}
\setotherlanguage{english}

\newfontfamily\hebrewfont{Magalith Tanach Pro}[Path = ../../vendor/fonts/]
% Last-resort fit for an overflowing page (inserted by build-body.mjs at the
% top of the body): the same font squeezed horizontally, e.g. 0.98 = letters 2% narrower.
\newcommand{\condensebody}[1]{%
  \renewfontfamily\hebrewfont{Magalith Tanach Pro}[Path = ../../vendor/fonts/, FakeStretch = #1]%
  \hebrewfont\fontsize{27}{30.6pt}\selectfont}
% Fit step after that (build-body.mjs STRICT_LEADING): lines exactly one
% \baselineskip apart, even where tall accent stacks would make TeX add \lineskip.
\newcommand{\strictleading}{\lineskiplimit=-\maxdimen}

% Every verse/chapter number is a \marginpar, i.e. a float. LaTeX's default 18
% float slots ran out on pages with many short verses (Psalms, Job) and the
% extra numbers were silently dropped (page 1003, Ps 22:27-31).
\extrafloats{100}
% Verse number size in pt (the scan's are ~17pt; was 9).
\newcommand{\vnumsize}{14}
% Margin number set against the text: the margin box is right-aligned (Hebrew),
% which suits odd pages (left margin) but put even-page numbers at the far side
% of the box, and let a number wider than the box (34pt gematria קכא, page
% 1072) spill into the text. Anchored on the text side, it spills outward.
\newcommand{\margintotext}[1]{\hbox to \linewidth{\ifthenelse{\equal{\pageparity}{even}}{\hss#1}{#1\hss}}}
\newcommand{\vnum}[1]{\leavevmode\marginpar{\margintotext{\textenglish{\fontsize{\vnumsize}{\vnumsize}\selectfont #1}}}}

% Chapter heading on its own centered line, body-size text. The text line
% before it ends ragged (like a petuchah): the body otherwise forces every
% line full width (\parfillskip=0pt), which spread a chapter's short last
% line edge to edge (page 74).
% Optional #1 = the label before the number (default "Chap."; Psalms: "Psalm"),
% set in small caps - a full capital, then smaller capitals - like the scan's
% "Cap." / "Psal." (\chaplabel, also used in the running head).
\newcommand{\chaplabel}[1]{\textsc{#1}}
\newcommand{\chapheading}[3][Chap.]{{\parfillskip=0pt plus 1fil\par}{\centering\fontsize{27}{30.6pt}\selectfont{\hebrewfont #3}\ \textenglish{\chaplabel{#1}~#2}\par}}
% Smaller chapter heading (20pt) for pages where space is tight
\newcommand{\chapheadingsmall}[3][Chap.]{{\parfillskip=0pt plus 1fil\par}{\centering\fontsize{20}{22pt}\selectfont{\hebrewfont #3}\ \textenglish{\chaplabel{#1}~#2}\par}}

% Chapter marker in margin, title-size. A gematria wider than \chapnumroom
% (קכא at 34pt, page 1072) would run off the paper: the chapter's first line
% gives up the excess on the margin's side (even pages: the line's start is
% nudged in; odd pages: it ends early) and the number moves into that gap.
% Outer margin 47pt minus \marginparsep: every 1-2 letter gematria fits (widest
% מה 34.3pt); 3 letters (קיט 40.7pt .. קמח 50.3pt) nudge the first line.
\newcommand{\chapnumroom}{36pt}
\newdimen\chapnumexcess
\newcommand{\chapnum}[1]{\leavevmode
  \setbox0\hbox{\hebrewfont\fontsize{34}{34}\selectfont #1}%
  \chapnumexcess=\dimexpr\wd0-\chapnumroom\relax
  \ifdim\chapnumexcess>0pt
    \ifthenelse{\equal{\pageparity}{even}}%
      {\parshape 2 0pt \dimexpr\linewidth-\chapnumexcess\relax 0pt \linewidth}%
      {\parshape 2 \chapnumexcess \dimexpr\linewidth-\chapnumexcess\relax 0pt \linewidth}%
  \else \chapnumexcess=0pt \fi
  \marginpar{\ifthenelse{\equal{\pageparity}{even}}%
    {\hbox to \linewidth{\hss\copy0\kern-\chapnumexcess}}%
    {\hbox to \linewidth{\kern-\chapnumexcess\copy0\hss}}}}

% פ/ס at body size, each centred in the gap it makes (as in the scan).
% The letters are printed only in the Torah: in the Prophets and Writings the
% scan keeps the same gaps and open lines but leaves the letter out (Josh 1:9,
% page 352; Job 2:10, page 1120), and a noted break keeps its masora circle
% alone in the gap (Josh 4:1, page 356). build-body.mjs starts non-Torah pages
% with \breaklettersfalse.
\newif\ifbreakletters \breakletterstrue
\newcommand{\breakletter}[2]{\ifbreakletters{#1#2}\else{\hebrewfont #2}\fi}
% Petuchah (open gap): the rest of the line is left open, the פ in its middle.
% Optional argument: marks printed on the letter (the masora circle when noted).
% Tied to the preceding word (\unskip\nobreak) so the פ never sits alone on a
% line; \nobreak + \null keep the closing half of the gap on the same line.
\newcommand{\petuchah}[1][]{\openbreakletter{\breakletter{פ}{#1}}}
% A setumah that ends a paragraph (before a chapter heading, portion end, ...):
% build-body.mjs uses this instead of \setumah, so it too sits mid-gap.
\newcommand{\setumahopen}[1][]{\openbreakletter{\breakletter{ס}{#1}}}
% Both halves of the gap stretch equally; only the first has a minimum (a word
% space), so on a full line the letter still ends the line, as it did before.
\newcommand{\openbreakletter}[1]{\unskip\nobreak\hskip\fontdimen2\font plus 1fil{#1}\nobreak\hskip 0pt plus 1fil\null{\parfillskip=0pt\par}}
% Setumah (closed gap): mid-line, or ending a line as in the scan (page 227,
% Num 6:23 and 6:26) - never starting one. The line may break only after the ס:
% the word space before it is replaced by an unbreakable one (\unskip\nobreak\ ),
% and the gap after it follows the break point, so at a line end it is dropped
% and the ס ends the line. The gap is 3em with the ס in its middle, and may
% shrink to 1em, as in the scan, so a short word after it isn't pushed alone
% onto a new line (page 121).
\newcommand{\setumah}[1][]{\unskip\nobreak\ \hspace{1.5em minus 1em}\breakletter{ס}{#1}\penalty0\hspace{1.5em minus 1em}}
% Setumah on a line widened by build-body.mjs (wide letters): fixed gap #2 on
% each side, so the line's words take the space instead (page 227).
\newcommand{\setumahtight}[2][]{\unskip\nobreak\ \hspace{#2}\breakletter{ס}{#1}\penalty0\hspace{#2}}
% Setumah as the last thing on a page: nothing follows, so no gap - just
% the ס with a space before it.
\newcommand{\setumahpageend}[1][]{{} \breakletter{ס}{#1}}
\newcommand{\parashabreak}{\par\vspace{4\baselineskip}}
% Masora circle of an unrendered book marker (װ), standing alone just before
% the book's first word.
\newcommand{\masoracircle}{{\hebrewfont ֯}}

% End of a weekly portion: "פ פ פ" (open) or "ס ס ס" (closed) on its own
% centered line at chapter-heading size (small variant when the page is tight).
% #2 = פ or ס; optional #1 = marks on the first letter (masora circle).
% The text line before it ends ragged, like a petuchah.
\newcommand{\sedrahend}[2][]{{\parfillskip=0pt plus 1fil\par}{\centering\fontsize{27}{30.6pt}\selectfont{\hebrewfont #2#1\ #2\ #2}\par}}
\newcommand{\sedrahendsmall}[2][]{{\parfillskip=0pt plus 1fil\par}{\centering\fontsize{20}{22pt}\selectfont{\hebrewfont #2#1\ #2\ #2}\par}}
% Last resort when even the small line overflows: inline at the end of the
% portion's last text line, body size, no extra line. Open (פ פ פ) ends the
% line like a petuchah; closed (ס ס ס) keeps the setumah gap and the text continues.
\newcommand{\sedrahendinlineopen}[1][]{\unskip\nobreak\ {פ#1\ פ\ פ}{\parfillskip=0pt plus 1fil\par}}
\newcommand{\sedrahendinlineclosed}[1][]{{} {ס#1\ ס\ ס}\hspace{3em}\penalty10000}

% Book title for a book that begins mid-page (page 499, 2 Samuel): the previous
% book's last line ends ragged, then a short rule, Hebrew name over English
% name, a short rule, and the text. Small variant when the page is tight.
% #1 = Hebrew name, #2 = English name. Sizes / gaps:
\newcommand{\midtitlerule}{{\centering\rule{30mm}{1pt}\par}}
\newcommand{\midtitlebefore}{4mm}     % previous text -> rule
\newcommand{\midtitlerulegap}{4mm}    % rule <-> names
\newcommand{\midtitlegap}{2mm}        % between the two names
\newcommand{\midtitleafter}{4mm}      % rule -> text
\newcommand{\midbooktitlebody}[6]{{\parfillskip=0pt plus 1fil\par}%
  \vspace{#5}\midtitlerule\vspace{#6}%
  {\centering{\fontsize{#3}{#3}\selectfont #1\par}\vspace{\midtitlegap}%
   {\fontsize{#4}{#4}\selectfont\textenglish{#2}\par}}%
  \vspace{#6}\midtitlerule\vspace{\midtitleafter}}
\newcommand{\midbooktitle}[2]{\midbooktitlebody{#1}{#2}{50pt}{30pt}{\midtitlebefore}{\midtitlerulegap}}
\newcommand{\midbooktitlesmall}[2]{\midbooktitlebody{#1}{#2}{36pt}{22pt}{1mm}{2mm}}

% Sub-book title (the five books of Psalms, work/psalms-books.json), before the
% sub-book's first chapter heading, as in the scan (page 1016, Ps 42): Hebrew
% "ספר שני" over English "BOOK 2" (the scan's Latin "LIBER SECUNDUS"). The
% previous line ends ragged. #1 = Hebrew, #2 = English. Sizes / gaps:
\newcommand{\subtitlebefore}{4mm}     % previous text -> Hebrew
\newcommand{\subtitlegap}{1mm}        % between the two lines
\newcommand{\subtitleafter}{2mm}      % English -> chapter heading
\newcommand{\subbooktitlebody}[5]{{\parfillskip=0pt plus 1fil\par}%
  \vspace{#5}%
  {\centering{\fontsize{#3}{#3}\selectfont{\hebrewfont #1}\par}\vspace{\subtitlegap}%
   {\fontsize{#4}{#4}\selectfont\textenglish{#2}\par}}%
  \vspace{\subtitleafter}}
\newcommand{\subbooktitle}[2]{\subbooktitlebody{#1}{#2}{30pt}{20pt}{\subtitlebefore}}
\newcommand{\subbooktitlesmall}[2]{\subbooktitlebody{#1}{#2}{22pt}{15pt}{1mm}}

% End of a book whose successor starts on a new page: "חזק" in place of the
% "פ פ פ" line, same sizes; the inline variant ends the last text line.
\newcommand{\chazak}{{\parfillskip=0pt plus 1fil\par}{\centering\fontsize{27}{30.6pt}\selectfont{\hebrewfont חזק}\par}}
\newcommand{\chazaksmall}{{\parfillskip=0pt plus 1fil\par}{\centering\fontsize{20}{22pt}\selectfont{\hebrewfont חזק}\par}}
\newcommand{\chazakinline}{\unskip\nobreak\ {חזק}{\parfillskip=0pt plus 1fil\par}}

% Masora finalis (word_notes.json reason 15), after a book's last verse, above
% the footer (the page's masora notes stay in the footer). #1 = the summary
% lines. Each \finalisline is one scan line: a long line is justified to the
% full width, a short one is centered. \finalisbig = words set in large type
% (**...** in the note).
\newlength{\finaliswd}
\newcommand{\finalissize}{\fontsize{22}{30.6pt}\selectfont}
\newcommand{\finalisbig}[1]{{\fontsize{28}{30.6pt}\selectfont #1}}
\newcommand{\finalisfullline}{0.8\linewidth}
\newcommand{\finalisline}[1]{{\finalissize
  \settowidth{\finaliswd}{\beginR #1\endR}%
  \ifdim\finaliswd>\finalisfullline \hbox to \linewidth{\beginR #1\endR}%
  \else \hbox to \linewidth{\hfil\beginR #1\endR\hfil}\fi}}
% A חזק line other than plain חזק (חזק ונתחזק, ...), at חזק size.
\newcommand{\finalischazak}[1]{{\centering\fontsize{27}{30.6pt}\selectfont{\hebrewfont #1}\par}}
\newcommand{\finalisblock}[1]{{\parfillskip=0pt plus 1fil\par}%
  \vspace{3mm}#1\par
  \vspace{3mm}{\centering\rule{4cm}{1.2pt}\par}\vspace{3mm}}

% Songs (Exod 15, ...), laid out line by line from the scan by build-body.mjs
% (work/song-layouts.json). The text before a song ends ragged; each scan line
% is forced (\songbreak) and its blocks pushed apart (\songgap), so 2 blocks
% sit flush right / flush left and the middle one of 3 between them. A
% one-block line is justified like any other body line.
\newcommand{\songstart}{\ifhmode{\parfillskip=0pt plus 1fil\par}\fi}
\newcommand{\songbreak}{\unskip\break}
\newcommand{\songgap}{\unskip\hfill}
% An empty first block: the line's only block goes flush left (Esth 9:7, page 1199).
\newcommand{\songemptyblock}{\mbox{}\hfill}
% Two-column songs (Eccl 3:2-8): the second block sits in a box whose right
% edge is #1 of the line width from the left edge, so every line's second
% block starts at the same column, as in the scan (page 1177).
\newcommand{\songcolumn}[1]{\unskip\hfill\hbox to #1\linewidth\bgroup}
\newcommand{\songcolumnend}{\unskip\hfill\egroup}

% A weekly portion that ends with no פ/ס mid-book (Gen47:27): the line ends
% ragged and the next portion starts on a new line, with no marker printed.
\newcommand{\portionbreak}{{\parfillskip=0pt plus 1fil\par}}

% Line filling with wide letterforms (jalt), chosen letter by letter by
% build-body.mjs for lines whose interword space would otherwise be too large.
% Applied to one letter (+ its marks) at a time: jalt widens every covered
% glyph it touches, so it must never span more than the chosen letter.
\newcommand{\wide}[1]{{\addfontfeature{RawFeature=+jalt}#1}}

% Measurement marks for that choice (build-body.mjs measure pass only): each
% word part is bracketed by \wb/\we, which record their page position at
% shipout into pageN.pos. \LR keeps "save position, then write" in that order
% (in right-to-left text the pair would otherwise be reversed and every write
% would report the previous mark's position).
\newwrite\posfile
\newcommand{\posopen}{\immediate\openout\posfile=page\pagecfg.pos}
\newcommand{\wb}[1]{\LR{\pdfsavepos\write\posfile{b #1 \the\pdflastxpos\space\the\pdflastypos}}}
\newcommand{\we}[1]{\LR{\pdfsavepos\write\posfile{e #1 \the\pdflastxpos\space\the\pdflastypos}}}
% Body-font metrics for the same pass: normal space, a reference letter, and
% each wide-capable letter's normal and wide width (#1 = the letters).
\newcommand{\posmetrics}[1]{%
  \immediate\write\posfile{space \the\fontdimen2\font}%
  \setbox0\hbox{ב}\immediate\write\posfile{unit \the\wd0}%
  \posmetricsloop#1\relax}
\def\posmetricsloop#1{\ifx\relax#1\else
  \setbox0\hbox{#1}\setbox2\hbox{\wide{#1}}%
  \immediate\write\posfile{letter #1 \the\wd0 \space\the\wd2}% (the plain space ends the box number)
  \expandafter\posmetricsloop\fi}

% Tolerance settings.
\tolerance=9999
\emergencystretch=2em

% Sized per context.
\newcommand{\largeletter}[1]{{\hebrewfont\fontsize{34}{20}\selectfont #1}}
\newcommand{\smallletter}[1]{{\hebrewfont\fontsize{16}{20}\selectfont #1}}
% Hanging letter (נ׳ תלויה, Judg 18:30; ע׳ תלויה, Ps 80:14, Job 38:13,15): small
% size (\hangsize), raised so its top stands \hangabove (a fraction of the
% full-size letter height) above the tops of the full-size letters.
% \hangletter{letter}{letter + marks above}{marks below} (build-body.mjs
% hangLetterTex): the marks above rise with the small letter; the marks below
% stay full size on the line. They are set on \hangcarrier, a letter with a
% flat base, drawn clipped to below its depth so only the marks show (the
% letter's own outline can't be clipped: ע's hiriq tucks up into its tail,
% Job 38:13). The space is the full-size letter's width (or the small
% letter's, if wider), with the marks and the small letter centred in it.
\newcommand{\hangcarrier}{ב}
\newsavebox\hangbase \newsavebox\hangup \newsavebox\hanglow
\newdimen\hangraise \newdimen\hangwidth
\newcommand{\hangsize}{20}
\newcommand{\hangabove}{0.2}
\newcommand{\hangletter}[3]{\sbox\hangbase{#1}\sbox\hanglow{\hangcarrier#3}%
  \sbox\hangup{\hebrewfont\fontsize{\hangsize}{20}\selectfont #2}%
  \setbox0\hbox{ב}\setbox2\hbox{\hebrewfont\fontsize{\hangsize}{20}\selectfont ב}%
  \hangraise=\dimexpr\ht0-\ht2+\hangabove\ht0\relax
  \hangwidth=\ifdim\wd\hangbase>\wd\hangup \wd\hangbase\else\wd\hangup\fi
  \LR{\tikz[baseline=0pt]{%
  \useasboundingbox (0,-\dp\hanglow) rectangle (\hangwidth,\ht0);
  \node[inner sep=0,anchor=base west] at ({(\hangwidth-\wd\hangup)/2},\hangraise) {\usebox\hangup};
  \begin{scope}\clip (-1em,-2em) rectangle ([xshift=1em]\hangwidth,{-\dp0-0.01em});
    \node[inner sep=0,anchor=base west] at ({(\hangwidth-\wd\hanglow)/2},0) {\usebox\hanglow};\end{scope}}}}
% Cut letter (ו׳ קטיעה, Num 25:12): \cutletter{cut}{gap}{shift}{letter}. The
% glyph is drawn twice in the same place: one copy clipped to below `cut`
% (height above the baseline) and moved right by `shift`, the other clipped to
% above cut+gap - a gap in the stem with the lower piece offset. Marks on the
% letter (masora circle) stay with the upper piece; the width is unchanged.
% Values from build-body.mjs (CUT_LETTER_*).
\newsavebox\cutbox
\newcommand{\cutletter}[4]{\sbox\cutbox{#4}\LR{\tikz[baseline=0pt]{%
  \useasboundingbox (0,0) rectangle (\wd\cutbox,\ht\cutbox);
  \begin{scope}\clip (-1em,-1em) rectangle ([xshift=1em]\wd\cutbox,#1);
    \node[inner sep=0,anchor=base west] at (#3,0) {\usebox\cutbox};\end{scope}
  \begin{scope}\clip (-1em,{#1+#2}) rectangle ([xshift=1em]\wd\cutbox,2em);
    \node[inner sep=0,anchor=base west] at (0,0) {\usebox\cutbox};\end{scope}}}}

% Qere/Ketiv at 3/4 body size (20.25pt against 27pt).
\newcommand{\qerisub}[1]{\raisebox{-0.5ex}{{\hebrewfont\fontsize{20.25}{20.25}\selectfont #1}}}
\newcommand{\qerisup}[1]{\raisebox{0.8ex}{{\hebrewfont\fontsize{20.25}{20.25}\selectfont #1}}}

% Masora-note apparatus (notes flow horizontally on one line, separated by spaces).
\newcommand{\pagefootnotes}{}
% Each note is one unbreakable box (pointer + text); footer lines break only
% between notes - except a note too long for a line, which wraps inside itself
% like ordinary text (Josh 21:36, page 389).
\newsavebox{\masoranotebox}
\newcommand{\masoranotetext}[2]{{\fontsize{12}{14.5}\selectfont \textdegree~\textenglish{#1}~#2}}
\newcommand{\masoranote}[2]{\sbox{\masoranotebox}{\masoranotetext{#1}{#2}}%
  \ifdim\wd\masoranotebox>0.85\textwidth \masoranotetext{#1}{#2}\else\usebox{\masoranotebox}\fi\hspace{2em}}
% A note's **...** part (verses the scan prints under the notes, e.g. the two
% verses after Josh 21:35 found only in some books): build-body.mjs moves it to
% \pagefootverses, set in large type below a second rule, with an asterisk.
\newcommand*{\pagefootverses}{}  % * (not \long): page-master tests it with \ifx against \empty
\newcommand{\footverses}[1]{{\fontsize{22}{27}\selectfont\hebrewfont *#1\par}}

% Build chapter-list header: takes a comma-separated list of chapter numbers
% (arabic) and a corresponding comma-separated list of gematria letters,
% outputs the formatted header line for odd/even pages.
% Internal macro to format a single chapter pair (arabic + gematria).
\def\@chapitem#1#2{ #2 \ \textenglish{#1}}

% Build a header line from chapter lists (both comma-separated), odd or even parity.
% #1 = parity (odd/even), #2 = chapter numbers (comma-sep), #3 = gematrias (comma-sep),
% #4 = english book name
\newcommand{\buildcapheader}[4]{%
  \ifthenelse{\equal{#1}{odd}}{%
    % Odd: gematrias first (in visual RTL order so they read first),
    % then arabic numbers, then book name.
    {\hebrewfont #3}\ \textenglish{#2} \ \textenglish{Cap.} — \textenglish{#4}%
  }{%
    % Even: hebrew bookname first, then chapters.
    \textenglish{Cap.} \textenglish{#2} — {\hebrewfont #3}%
  }%
}

% Main setup macro: call this once per page with chapter/catchword data.
% #1 = odd|even, #2 = hebrew bookname (e.g., בראשית), #3 = english bookname (e.g., GENESIS)
% #4 = chapters (comma-sep arabic, e.g., 1,2), #5 = gematrias (comma-sep hebrew, e.g., א,ב)
% #6 = catchword for footer
\newcommand{\setuppage}[6]{%
  \ifthenelse{\equal{#1}{even}}{\reversemarginpar}{}%
  \fancypagestyle{bookpage}{%
    \fancyhf{}%
    \renewcommand{\headrulewidth}{0.4pt}%
    \ifthenelse{\equal{#1}{even}}{%
      \fancyhead[R]{\begin{hebrew}\textenglish{\fontsize{10}{10}\selectfont\thepage}\end{hebrew}}%
    }{%
      \fancyhead[L]{\begin{hebrew}\textenglish{\fontsize{10}{10}\selectfont\thepage}\end{hebrew}}%
    }%
    \fancyhead[C]{\begin{hebrew}{\fontsize{11}{11}\selectfont%
      \ifthenelse{\equal{#1}{odd}}{%
        {\hebrewfont #5}\ \textenglish{#4} \ \textenglish{Cap.} — \textenglish{#3}%
      }{%
        \textenglish{#2} — \textenglish{Cap.} \textenglish{#4} \ {\hebrewfont #5}%
      }%
    }\end{hebrew}}%
    \fancyfoot[C]{\begin{hebrew}%
      \rule{\textwidth}{0.4pt}\\[0.5mm]%
      \noindent\begin{tabular*}{\textwidth}{@{}p{0.5\textwidth}@{\extracolsep{\fill}}l@{}}
        \raggedright{\fontsize{10}{12}\selectfont\pagefootnotes} & {\fontsize{10}{12}\selectfont #6}%
      \end{tabular*}%
    \end{hebrew}}%
  }%
  \pagestyle{bookpage}%
  \setlength{\parindent}{0pt}%
  \setlength{\parskip}{1mm}%
}

\providecommand{\pageparasha}{}
\providecommand{\pagebookstart}{false}\providecommand{\pagechaplabel}{Chap.}\newlength{\catchwidth}

% Book-opening title (pages with \pagebookstart{true}): printed at the top of
% the text block, Hebrew name over English name. Sizes / gaps:
% Room is limited by page 1 (least slack of the book-opening pages).
\newcommand{\booktitlehebsize}{\fontsize{58}{64}\selectfont}
\newcommand{\booktitleengsize}{\fontsize{32}{38}\selectfont}
\newcommand{\booktitlebefore}{2mm}    % between the header and the title
\newcommand{\booktitlegap}{4mm}       % between the two names
\newcommand{\booktitleafter}{10mm}    % between the title and the text

% Build the header line for odd vs. even parity.
% Odd: chapters first, then english bookname.
% Even: hebrew bookname first, then chapters.
\fancypagestyle{bookpage}{%
  \fancyhf{}%
  \renewcommand{\headrulewidth}{0.4pt}%
  \ifthenelse{\equal{\pageparity}{even}}{%
    \fancyhead[R]{\begin{hebrew}\textenglish{\fontsize{10}{10}\selectfont\pagecfg}\end{hebrew}}%
  }{%
    \fancyhead[L]{\begin{hebrew}\textenglish{\fontsize{10}{10}\selectfont\pagecfg}\end{hebrew}}%
  }%
  \ifthenelse{\equal{\pagebookstart}{true}}{%
    % Book-opening page: page number only, no rule, no chapter indicator -
    % the title itself is printed at the top of the text block (below).
    \renewcommand{\headrulewidth}{0pt}%
  }{\ifthenelse{\equal{\pageparity}{odd}}{%
    % Odd: gematrias א ב ... + chapters 1 2 ... + Chap. — BOOKNAME
    \fancyhead[C]{\begin{hebrew}{\fontsize{11}{11}\selectfont%
      {\hebrewfont\pagegematrias}\ \textenglish{\pagechapters}\ifthenelse{\equal{\pagechaplabel}{}}{}{\ \textenglish{\chaplabel{\pagechaplabel}}}\ —\ \textenglish{\pageengbook}%
    }\end{hebrew}}%
  }{%
    % Even: HEBBOOK — [PARASHA —] gematrias Chap.chapters (page 2: בראשית — א Chap.1).
    % \pageparasha is set only on Torah pages (from PARSHIOT.json); empty = omitted.
    \fancyhead[C]{\begin{hebrew}{\fontsize{11}{11}\selectfont%
      \pagehebbook\ —\ %
      \ifthenelse{\equal{\pageparasha}{}}{}{\pageparasha\ —\ }%
      \pagegematrias\ \textenglish{\ifthenelse{\equal{\pagechaplabel}{}}{}{\chaplabel{\pagechaplabel}}\pagechapters}%
    }\end{hebrew}}%
  }}%
  % Footer: the catchword sits at the far left of the first line (zero-height
  % overlay); the notes' first line stops 1em short of it, later lines use the
  % full width (\parshape). Each note is unbreakable unless longer than a line
  % (\masoranote). Large-type verses from a note (\pagefootverses) follow under
  % a second rule, as in the scan (page 389).
  \fancyfoot[C]{\begin{hebrew}%
    \settowidth{\catchwidth}{\fontsize{14.5}{14.5}\selectfont\pagecatchword}%
    \noindent\hbox to \textwidth{\rule{\footerrulewidth}{0.4pt}\hfill}\\[2mm]%
    \begin{minipage}[t]{\textwidth}%
      \setlength{\parindent}{0pt}%
      \vbox to 0pt{\hbox to \textwidth{\hfill{\fontsize{14.5}{14.5}\selectfont\pagecatchword}}\vss}%
      \raggedright
      \parshape 2 0pt \dimexpr\textwidth-\catchwidth-1em\relax 0pt \textwidth
      \pagefootnotes\par
      \ifx\pagefootverses\empty\else
        \vspace{1mm}\noindent\rule{\textwidth}{0.4pt}\par\vspace{1mm}%
        \parshape 0 \raggedright\pagefootverses
      \fi
    \end{minipage}%
  \end{hebrew}}%
}
% \pagestyle{bookpage} picks the header layout from the page's settings when it
% runs, so page-content.tex selects it per page (after the settings).
\flushbottom
% Max stretch per line gap, in pt (2 = 30.6pt -> 32.6pt; fills 456 of 495
% built pages completely).
\newcommand{\bodylinestretch}{2}
\newdimen\bodyslack \newdimen\bodystretchcap
% Tallest notes block that stays in place under the rule; anything taller moves
% the rule up (calibrated: 2 rows fit with >= 0.5in to the paper edge).
\newdimen\footnotesroom \footnotesroom=40pt
% Rule above the footer: starts at the right (Hebrew) edge of the text, as in
% the scan (~0.4 of the text width; measured on pages 902, 919, 1043, 1200).
\newcommand{\footerrulewidth}{0.4\textwidth}
\setlength{\parindent}{0pt}
\setlength{\parskip}{1mm}

% Per-page state that a page may change and the next page must get back (one-run
% book): the text block height (moved up for tall notes / footer verses) and the
% Hebrew font (narrowed by \condensebody).
\newdimen\booktextheight \booktextheight=\textheight
\let\bookhebrewfont\hebrewfont
\makeatletter
\newcommand{\bookpagereset}{%
  \global\textheight=\booktextheight
  \global\vsize=\textheight
  \global\@colht=\textheight\global\@colroom=\textheight
  \let\hebrewfont\bookhebrewfont}
\makeatother
